\chapter{What Git Actually Does}
\label{ch:basics}

\chapabstract{The minimum theory needed before the commands make sense:
what a repository is, what a commit is, and where your files actually
live at each stage.}

\section{Version control, in one sentence}

\Git{} is a tool that records snapshots of a set of files over time, so
that you can go back to any earlier snapshot, see what changed between
two snapshots, and combine changes made by different people. Each
snapshot is called a \textbf{commit}. Unlike copying a folder by hand,
every commit is linked to the one before it, forming a history you can
inspect, search and revert.

\Git{} is \textbf{distributed}: every clone of a repository is a full
copy of the entire history, not just a checkout of the current state.
There is no central server that must be reachable to commit, browse
history, or create a branch; a server is only needed to synchronise with
other people.

\section{Three places your files can be}

Understanding \Git{} means understanding that a tracked file exists in
up to three places at once.

\begin{enumerate}
  \item \textbf{Working directory.} The files you see and edit on disk.
  \item \textbf{Staging area} (also called the \emph{index}). A
        holding area where you assemble exactly the changes that will
        go into the next commit. This is what makes \Git{} different
        from most other tools: you choose what to commit, not
        everything you touched.
  \item \textbf{Repository.} The database of committed snapshots,
        stored in the hidden \file{.git} folder. Once a change is
        committed, it is part of the permanent history.
\end{enumerate}

\begin{figure}[H]
  \centering
  \begin{tikzpicture}[
      node distance=1.4cm,
      box/.style={draw=codekey,fill=codebg,line width=1pt,
                  rounded corners=3pt,minimum width=3.1cm,
                  minimum height=1.3cm,align=center,font=\sffamily\small},
      lbl/.style={font=\scriptsize\sffamily,midway,above}
    ]
    \node[box] (wd) {Working\\Directory};
    \node[box,right=of wd] (idx) {Staging Area\\(Index)};
    \node[box,draw=accentalt,right=of idx] (repo) {Local\\Repository};
    \draw[-{Latex[length=3mm]},line width=1.1pt,codekey]
      (wd.east) -- (idx.west) node[lbl]{\code{git add}};
    \draw[-{Latex[length=3mm]},line width=1.1pt,accentalt]
      (idx.east) -- (repo.west) node[lbl]{\code{git commit}};
    \draw[-{Latex[length=3mm]},line width=1.1pt,codenumber]
      (repo.south) --++(0,-8mm) -| (wd.south)
      node[pos=0.25,below,font=\scriptsize\sffamily]{\code{git checkout} / \code{git restore}};
  \end{tikzpicture}
  \caption{The everyday local workflow: edit, stage, commit.}
  \label{fig:basics:workflow}
\end{figure}

\begin{gotcha}{"I saved the file, so it's in Git"}
Saving a file only changes the working directory. It is not staged
until you run \code{git add}, and it is not part of the history until
you run \code{git commit}. Three separate, deliberate steps.
\end{gotcha}

\section{Commits, hashes and branches}

A commit is a snapshot of the whole project (not a diff), together with
an author, a timestamp, a message, and a pointer to its parent commit.
\Git{} identifies each commit by a SHA-1 hash (\eg
\code{a3f1c9e...}), computed from its content: change one byte anywhere
in the project and every commit made afterwards gets a new hash.

A \textbf{branch} is nothing more than a movable pointer to a commit.
The default branch is conventionally called \code{main} (older
repositories use \code{master}). Creating a branch is instantaneous and
cheap because it does not copy any files, only records a new pointer.
\code{HEAD} is a pointer to whichever branch (or commit) you currently
have checked out.

\begin{ruleofthumb}
If you remember one thing from this chapter: a commit is a snapshot, a
branch is a label on a commit, and \code{HEAD} says which label you are
standing on.
\end{ruleofthumb}

\section{Local versus remote}

Everything so far happens on your machine. A \textbf{remote} is simply
another copy of the repository, usually hosted on a server such as
\GitHub{}, identified by a short name (conventionally \code{origin}) and
a URL. \code{git push} sends your local commits to a remote branch;
\code{git fetch} downloads commits from a remote without touching your
working directory; \code{git pull} does a fetch followed by a merge (or
rebase) into your current branch. \Cref{ch:github} covers this in
detail.

\section{Checklist}
\label{sec:basics:checklist}

\begin{itemize}
  \item A file only enters history through \textbf{add} then
        \textbf{commit}; there is no implicit save.
  \item A commit is a full snapshot identified by a content hash, with
        a pointer to its parent.
  \item A branch is a pointer, not a copy; \code{HEAD} shows which one
        you are on.
  \item Every clone holds the entire history; a remote is just another
        copy to synchronise with.
\end{itemize}
